\section{Conclusão}

Este trabalho apresenta uma arquitetura multiagente baseada em LLMs para
diagnóstico personalizado de carteiras financeiras. A proposta combina
agentes especializados, debate adversarial, moderação estruturada e
roteamento heterogêneo de modelos, mantendo a saída limitada a
diagnóstico e evitando recomendações de compra ou venda.

A contribuição central está na avaliação. O sistema é comparado contra
baselines fortes, incluindo agente único com self-consistency e
multiagente sem debate, e isso permite separar ganhos reais de
personalização e cobertura de risco de melhorias meramente textuais.
Os resultados mostram que a validação estrutural, sem dados externos,
expõe melhor como os sistemas multiagente tratam lacunas de evidência. A
comparação principal, feita com 20 carteiras enriquecidas, mostra diferenças
mais úteis entre os modos de execução: custo, latência, conclusividade,
detalhamento e concordância com consenso. Nesse cenário financeiro, B1 é o
mais barato, B2 é o mais detalhista em número de fatores, B3 apresenta o
melhor equilíbrio geral, e B4-R reduz custo e latência em relação a B4-H.

Assim, a principal conclusão não é que o debate adversarial já supera os
baselines em precisão financeira. Na execução enriquecida, B3 foi o melhor
equilíbrio entre concordância com consenso, diagnósticos conclusivos e custo
proxy, enquanto B4-R reduziu custo e latência em relação a B4-H, mas não
melhorou a proxy de acurácia. A conclusão mais defensável é que a
arquitetura proposta permite avaliar essa hipótese de forma controlada,
rastreável e comparável, e que o benefício do debate ainda precisa ser
demonstrado por métricas mais fortes. Como próximos passos, permanecem a
avaliação contrafactual por perfil, a verificação factual automática, a
calibração de LLM-as-judge com anotação humana e a comparação com desempenho
posterior das carteiras.
